\chapter{Fejlesztői dokumentáció}
\label{ch:impl}

\section{Követelményelemzés és specifikáció}

A fejlesztés és kód írás megkezdése előtt fontos egy rövid felmérést, kutatást
végeznünk és bizonyos szempontok alapján meg kell vizsgálnunk a feladatot.
A rendszer fejlesztésének első lépése a követelmények meghatározása és dokumentálása.
A specifikáció elkészítése.
 A követelményelemzés során tisztázzuk felhasználói igényeket, a rendszerrel szemben támasztott elvárásainkat,
valamint azokat a korlátokat és feltételeket, amelyek a fejlesztést befolyásolják.
 
   
A fejezet a specifikáció és követelménymenedzsment során alkalmazott
lépéseket követi. Ennek megfelelően először a rendszer megvalósíthatóságát
  vizsgálom, majd a követelmények feltárását és elemzését végzem el. 
Ezután bemutatom a funkcionális és nem funkcionális követelményeket, 
ismertetem a rendszer alap modelljeit, végül pedig felvázolom azokat a feltételeket, amelyeket,
ha a rendszer teljesít, akkor elfogadható.

%----------------------------------------------------------------------
\clearpage
\subsection{Megvalósíthatósági elemzés}

Itt vizsgálom a rendszer, a szoftver és a fejlesztés megvalósíthatóságát,
 figyelembe véve a rendelkezésre álló erőforrásokat, a költségeket és a rendelkezésre álló időtartamot,
  valamint az üzemeltetés és karbantartás kérdéseit.

A rendszer fejlesztését egyetlen, a tárgyat már elvégzett fejlesztő végzi, nincs külső erőforrás bevonására.
 A fejlesztés közvetlen költsége gyakorlatilag nulla, mivel a projekt kizárólag ingyenesen elérhető,
nyílt forráskódú technológiákra épül. (TypeScript, React, Vite, pnpm).
Hardverigény szempontjából a rendszer kedvező helyzetben van, mivel minden számítás a felhasználó
böngészőjében fut le, így nincs szükség külön szerver infrastruktúrára vagy felhőszolgáltatásra.
Az alkalmazás statikus weboldalként publikálható, például GitHub Pages szolgáltatáson keresztül,
amely ingyenes üzemeltetést tesz lehetővé.

A rendszer által végzett algoritmusok viszonylag kis méretű véges bemeneteken dolgoznak, ez által
 azok átlagos felhasználói eszközökön is gyorsan végrehajthatók.

A fejlesztés időkeretét egy egyetemi félév adja. Ennek megfelelően a funkciók prioritások szerint 
kerülnek megvalósításra, és az alacsonyabb prioritású elemek szükség esetén későbbi bővítésként kezelhetők.

A karbantartása egyszerű, ezt támogatja a, moduláris felépítés. 
Az üzleti logika elkülönül a felhasználói felülettől, a felület pedig a React komponensalapú felépítésére épül.
Így új funkciók hozzáadása vagy meglévő funkciók módosítása nem igényli a teljes rendszer
átírását. A karbantartást tovább segíti a Git alapú verziókezelés.

%----------------------------------------------------------------------
\subsection{Követelmények feltárása}
%----------------------------------------------------------------------

\subsubsection{Források és célcsoport}
A oldal tartalma, az előadások anyagára, órai jegyzetekre, valamit saját hallgatói tapasztalatra épül elsősorban.

 A célcsoport a Számításelmélet alapjai I--II. tárgy hallgatói, vizsgára készülők, 
 valamit az anyag mélyebb megértésére törekvő hallgatók.
\clearpage
\subsubsection{Hatókör}
% thesis.tex: "A rendszer célja és hatóköre" szöveg + tab:in_scope
% JAVÍTANDÓ: nincs grafikus automataszerkesztő; a táblázat oszlopdefiníciója {l X} legyen
% + ide jön a témakör–modul táblázat (tab:temakorok)
A szakdolgozat hatókörébe egy webalapú, oktatást támogató alkalmazás
tervezése és megvalósítása tartozik, amely a Számításelmélet
alapjai I--II. tárgy egyes témaköreinek megértését segíti interaktív
és vizuális eszközökkel. A~\ref{tab:in_scope}.~táblázat a hatókörbe
tartozó funkciókat foglalja össze.

\begin{table}[htbp]
      \centering 
      \begin{tabularx}{\textwidth}{|l|X|} 
            \hline 
            \textbf{Funkció} & \textbf{Leírás} \\ 
            \hline 
            Elméleti tartalomkezelés & 
                  Témakörönként strukturált, blokkalapú tananyag a Számításelmélet 
                  alapjai I--II. tananyagából: definíciók, tételek, lemmák, bizonyítások,
                  példák, képletek, táblázatok és figyelmeztetések stb. \\ \hline
            Interaktív algoritmusok &
                  A hallgató különböző bemenetek megadásával futtathatja a tananyagban szereplő algoritmusokat. 
                  A rendszer a végrehajtás eredményét és köztes állapotait is megjeleníti. \\ \hline
            Folyamatok vizualizációja &
                  Az alkalmazás grafikus elemek segítségével szemlélteti a számításelméleti folyamatokat, 
                  ezáltal támogatva az absztrakt fogalmak megértését. \\ \hline
            Önellenőrzés &
                  A rendszer gyakorló feladatokat és kvízeket biztosít az elsajátított ismeretek ellenőrzésére. \\ \hline
            Egységes tanulási környezet &
                  A rendszer egy helyen fogja össze a tananyagot, a vizualizációkat és az interaktív eszközöket, 
                  támogatva az önálló felkészülést. \\
            \hline
      \end{tabularx}
      \caption{A rendszer hatókörébe tartozó fő funkciók}
      \label{tab:in_scope}
\end{table}


\clearpage
\subsubsection{Hatókörön kívül eső elemek}
A hatókörön kívül eső funkciók a rendszer céljainak és 
erőforrásainak megfelelően kerültek meghatározásra. 
Ezek a funkciók nem részei a szakdolgozat keretein belül 
megvalósuló rendszernek, és a~\ref{tab:out_of_scope}.~táblázat foglalja össze őket.
% thesis.tex: tab:out_of_scope (változatlanul átvehető)
\begin{table}[H]
      \centering 
      \begin{tabularx}{\textwidth}{|>{\raggedright\arraybackslash}p{3.2cm}|>{\raggedright\arraybackslash}X|>{\raggedright\arraybackslash}X|}
            \hline 
            \textbf{Funkció} & \textbf{Leírás} & \textbf{Indoklás} \\ 
            \hline
            Teljes oktatási keretrendszer &
            Az alkalmazás nem helyettesíti az egyetemi oktatási
            rendszereket, az előadásokat, a gyakorlatokat vagy az oktatói
            magyarázatokat. A teljeskörű tananyag feldolgozás nem része és nem is célja a szakdolgozat keretein belül megvalósuló munkának.
            & Egy teljes oktatási keretrendszer kialakítása jelentősen túlmutat a szakdolgozat célkitűzésein és rendelkezésre álló erőforrásain. \\ \hline
            Felhasználókezelés &
            Az alkalmazás nem tartalmaz regisztrációs, bejelentkezési vagy
            profilkezelési funkciókat. & A rendszer működése nem igényel felhasználóspecifikus adatok tárolását vagy azonosítását. A felhasználókezelés bevezetése így nem nyújtana érdemi többletértéket, ugyanakkor növelné a rendszer komplexitását.  \\ \hline
            Mobilalkalmazás &
            A rendszer böngészőn keresztül használható, natív
            mobilalkalmazás nem készül.  & A rendszer böngészőalapú webalkalmazásként készül, így külön telepítés nélkül több platformon is használható. A natív mobilalkalmazás fejlesztése jelentős többletmunkát jelentene a szakdolgozat kereteihez képest.\\ \hline
            Valós idejű együttműködés &
            A rendszer nem támogatja több felhasználó egyidejű munkáját
            vagy a közös szerkesztést. & A rendszer elsődlegesen önálló tanulás támogatására készült. A többfelhasználós szinkronizációs platform kialalkítása, egy másik alapötlet és nem ez \\
            \hline
      \end{tabularx}
      \caption{A rendszer hatókörén kívül eső funkciók}
      \label{tab:out_of_scope}
\end{table}

\clearpage
\subsubsection{Prioritizálás}
% ea02, 5. dia: fontossági sorrend; minden elkészül, a prioritás a fejlesztési sorrendet adja
% - kötelező: témakörök + elméleti anyag, DVA-szimulátor
% - fontos: grammatika-levezetés, Chomsky-osztályozó, 3NF, regex-ellenőrző, kvíz, NDVA
% - utolsó kör: veremautomata, Turing-gép, logikai modul, localStorage
A fejlesztés során a funkciók az itt meghatározott prioritási sorrend
szerint készülnek el. A prioritás nem a funkciók fontosságát vagy
elhagyhatóságát jelöli, hanem a megvalósítás tervezett sorrendjét.

Elsőként az alkalmazás alapstruktúrája készül el: a navigáció, valamint
a tananyag tárolására szolgáló tartalomtípusok és a témakörök váza.
Ezek teszik lehetővé, hogy új tananyag a meglévő kód módosítása nélkül
legyen felvehető.

Ezt követően a fejlesztés témakörönként halad, a tárgy tematikus
sorrendjét követve. Minden témakörnél először az elméleti anyag kerül
feldolgozásra, majd az ehhez kapcsolódó algoritmus logikája és
interaktív, vizualizációs modulja. Így minden lépés végén egy működő,
kipróbálható állapot jön létre. Elsőként a Számításelmélet alapjai I.
tárgy témakörei készülnek el, mivel ezek alkotják a tárgy alapját,
és a II. rész is ezekre épül.

A témakörök és a hozzájuk tartozó modulok elkészültével a rendszer
eléri a kitűzött célokat: ez a változat önmagában is teljes értékű,
használható tanulási eszköz. Csak ezt követően kerülnek sorra a
tanulást kiegészítő funkciók, valamint az opcionális, továbbfejlesztett
képességek. Így ha a rendelkezésre álló idő nem elegendő mindenre,
a rendszer akkor is működőképes állapotban zárható le.

%----------------------------------------------------------------------
\subsection{Külső interfészek}
%----------------------------------------------------------------------
A rendszer nem kommunikál szerveroldali komponenssel,
adatbázissal vagy külső szolgáltatással, minden számítás a
böngészőben történik.
Külső interfészei ezért a~\ref{tab:interfaces}.~táblázatban felsoroltakra korlátozódnak.
\begin{table}[H]
    \centering
    \begin{tabularx}{\textwidth}{|p{4.2cm}|X|}
        \hline
        \textbf{Interfész} & \textbf{Leírás} \\
        \hline
        Felhasználói felület &
        A felhasználó webböngészőben, egérrel és billentyűzettel
        kezeli a rendszert. \\ \hline
        Böngésző tárhelye &
        A haladás és a beállítások a böngésző helyi tárolójában
        (localStorage) őrződnek meg. \\ \hline
        Kiszolgálás &
        Az alkalmazás statikus fájlokként kerül kiszolgálásra
        (pl.\ GitHub Pages). \\
        \hline
    \end{tabularx}
    \caption{A rendszer külső interfészei}
    \label{tab:interfaces}
\end{table}

%----------------------------------------------------------------------
\subsection{Funkcionális követelmények}
%--bevezesres kicsi szoveg
%----------------------------------------------------------------------
% TODO: FK-3 ("minden témakörhöz") a végén ellenőrizni
\begin{compactdesc}
    \item[FK-1] A rendszer biztosítsa a Számításelmélet alapjai I--II.
          tárgy témaköreinek elérését és kiválasztását, a tárgy
          tematikus sorrendjében.
    \item[FK-2] Minden témakörhöz tartozzon elméleti összefoglaló, amely
          a tárgyban használt jelöléseket és képleteket alkalmazza.
    \item[FK-3] Minden témakörhöz kapcsolódjon vizualizációs vagy
          interaktív modul vagy gyakorló feladat.
    \item[FK-4] A vizualizációs és interaktív modulok jelenítsék meg a
          számítási folyamatok és algoritmusok lépéseit, és tegyék
          lehetővé a lépések közötti előre- és visszalépést.
    \item[FK-5] A felhasználó adhasson meg saját bemenetet (szót,
          grammatikát, reguláris kifejezést, automatát) az ezt támogató
          modulok számára a rendszer dolgozza ezeket fel, és az eredményt
          egyértelműen jelenítse meg.
    \item[FK-6] Hibás vagy nem értelmezhető bemenet esetén a rendszer
          jelezze a hiba okát, és a modul maradjon használható.
    \item[FK-7] A rendszer biztosítson önellenőrző feladatokat (kvízeket),
          amelyekre a felhasználó azonnali visszajelzést kap.
    \item[FK-8] A rendszer a felhasználó haladását és beállításait a
          böngészőben helyben megőrizze.
\end{compactdesc}


A rendszer elsődleges szereplője a \textbf{Hallgató}, aki a webes
felületen keresztül éri el a tárgy tanulását támogató felületet. A
hallgató kiválaszthatja a számára érdekes témakört, megtekintheti az
elméleti összefoglalót, használhatja a vizualizációkat, és interaktív
feladatokon keresztül gyakorolhatja a tanultakat.

A vizualizációk segítségével a hallgató lépésről lépésre követheti az
egyes algoritmusok és modellek működését. Az interaktív feladatoknál
saját bemenetet adhat meg, amelyre a rendszer automatikusan
kiértékelést és visszajelzést ad.

A~\ref{fig:usecase}.~ábra a rendszer fő használati eseteit és azok
kapcsolatait mutatja be. Az \emph{«include»} kapcsolat azt jelöli, hogy
egy eset mindig magában foglal egy másikat: a vizualizáció és az
interaktív feladat is saját bemenettel indul, a feladat és a kvíz pedig
mindig visszajelzéssel zárul. Az \emph{«extend»} kapcsolat egy csak
bizonyos feltétel mellett bekövetkező esetet jelöl: hibás bemenetnél a
rendszer hibaüzenetet jelenít meg. A \emph{«precedes»} kapcsolat
előfeltételt fejez ki: a témakör kiválasztása megelőzi az elméleti anyag
és a hozzá tartozó modulok használatát. Az \emph{«invokes»} kapcsolat
rákövetkezést jelöl: egy témakör megnyitása, a kvíz kitöltése és a
nyelv választása után a rendszer automatikusan elmenti a haladást,
illetve a beállításokat a böngésző helyi tárolójába. A megnyitott
témaköröket a navigációs menü pipával jelöli, és kiírja, hogy a
félév témaköreiből eddig hányat nyitott meg a hallgató.

\begin{figure}[H]
    \centering
    \resizebox{\textwidth}{!}{\begin{tikzpicture}[
    wf,
    uc/.style={draw=black!70, ellipse, fill=black!4, minimum width=3.9cm,
               minimum height=1.25cm, align=center, inner sep=1pt},
    rel/.style={->, >=stealth, densely dashed, black!70},
    lbl/.style={font=\sffamily\scriptsize, fill=white, inner sep=1pt},
  ]
  \draw[black!70, line width=0.8pt] (2.2,-1.0) rectangle (14.6,11.4);
  \node[anchor=north west, font=\sffamily\footnotesize\bfseries] at (2.3,11.35) {Webes felület};

  \begin{scope}[shift={(0.8,4.4)}, black!80, line width=0.8pt]
    \draw (0,1.05) circle (0.25);
    \draw (0,0.8) -- (0,0.05);
    \draw (-0.45,0.55) -- (0.45,0.55);
    \draw (0,0.05) -- (-0.35,-0.5);
    \draw (0,0.05) -- (0.35,-0.5);
    \node[anchor=north] at (0,-0.55) {Hallgató};
  \end{scope}
  \coordinate (actor) at (1.25,4.95);

  \node[uc] (tema)  at (5.5,10.0) {Témakör\\kiválasztása};
  \node[uc] (elm)   at (5.5,8.0)  {Elméleti anyag\\megtekintése};
  \node[uc] (viz)   at (5.5,6.0)  {Vizualizáció\\futtatása};
  \node[uc] (fel)   at (5.5,4.0)  {Interaktív feladat\\megoldása};
  \node[uc] (kviz)  at (5.5,2.0)  {Kvíz kitöltése};
  \node[uc] (nyelv) at (5.5,0.0)  {Nyelv választása\\(HU / EN)};

  \node[uc] (hiba)  at (11.8,8.4) {Hibaüzenet\\megtekintése};
  \node[uc] (bem)   at (11.8,5.4) {Saját bemenet\\megadása};
  \node[uc] (vissz) at (11.8,2.4) {Visszajelzés\\megtekintése};
  \node[uc] (ment)  at (11.8,-0.2) {Haladás és beállítások\\mentése};

  \foreach \u in {tema, elm, viz, fel, kviz, nyelv}
    \draw[black!70] (actor) -- (\u.west);

  \draw[rel] (tema) -- node[lbl, right=2pt] {«precedes»} (elm);
  \draw[rel] (viz)  -- node[lbl, sloped] {«include»} (bem);
  \draw[rel] (fel)  -- node[lbl, sloped] {«include»} (bem);
  \draw[rel] (fel)  -- node[lbl, sloped] {«include»} (vissz);
  \draw[rel] (kviz) -- node[lbl, sloped] {«include»} (vissz);
  \draw[rel] (hiba) -- node[lbl, right] {«extend»} (bem);
  \draw[rel] (kviz) -- node[lbl, sloped] {«invokes»} (ment);
  \draw[rel] (nyelv) -- node[lbl, sloped] {«invokes»} (ment);
  \draw[rel] (tema.east) -- node[lbl, pos=0.5] {«invokes»} (14.2,10.0) -- (14.2,-0.2) -- (ment.east);
\end{tikzpicture}
}
    \caption{A rendszer használati esetei}
    \label{fig:usecase}
\end{figure}

\subsection{Nem funkcionális követelmények}
%----------------------------------------------------------------------
% ea02, 7-11. dia kategóriái szerint lehetne csoportosítani
% TODO: NFK-3 reszponzivitás ellenőrizni

\begin{compactdesc} \item[NFK-1 (Teljesítmény)] A felhasználói bemenetek kiértékelése és a vizualizációk egy lépésének kiszámítása átlagos asztali környezetben 1 másodpercen belül történjen meg, a potenciálisan nem termináló
számítások (pl. Turing-gép futtatása, szavak felsorolása)
lépés- vagy hosszkorláttal fussanak. 
\item[NFK-2 (Megbízhatóság)] 
Hibás vagy nem értelmezhető bemenet esetén a rendszer ne 
álljon le, hanem egyértelmű hibaüzenetet jelenítsen meg, 
miközben a funkcionalitás továbbra is használható marad. 
\item[NFK-3 (Hordozhatóság)] 
Az alkalmazás telepítés nélkül fusson a modern böngészők 
(Google Chrome, Microsoft Edge és Mozilla Firefox) aktuális 
verzióin. 
\item[NFK-4 (Használhatóság)] 
A rendszer felülete magyar nyelvű legyen, és a Számításelmélet
 alapjai I--II. tárgyban használt jelölésrendszert alkalmazza.
  \item[NFK-5 (Karbantarthatóság)]
  Az algoritmusok megvalósítása elkülönüljön a felhasználói
   felülettől, lehetővé téve új tananyagok és modulok 
   hozzáadását a meglévő rendszer jelentős módosítása nélkül. 
\end{compactdesc}


%----------------------------------------------------------------------
\subsection{Rendszermodellek}
%----------------------------------------------------------------------
% ea02, 6. és 42. dia: felület terv, adatformátum leírás

\subsubsection{A felhasználói felület modellje}
A felhasználói felület tervezésénél a fő szempont az áttekinthetőség
és a könnyű navigáció volt. Az alkalmazás két oldaltípusból áll: a
főoldalból és a témakörök oldalaiból. Mindkettő tetején ugyanaz a
fejléc látható, amelyben a két félév témakörei egy-egy lenyíló
menüből érhetők el.
\clearpage
A főoldal (\ref{fig:homepage}.~ábra) részei: (1) fejléc a
navigációval, (2) rövid bemutatkozó szöveg, (3) lábléc, (4) a
fejlécből lenyíló, tematikus sorrendű témakörlista.

\begin{figure}[htbp]
    \centering
    \resizebox{\textwidth}{!}{\begin{tikzpicture}[wf]
  \draw[wf frame] (0,0) rectangle (10,11);

  \draw[wf bar] (0,10.1) rectangle (10,11);
  \draw[black!70, line width=0.7pt] (0.35,10.35) -- (0.35,10.65) -- (0.55,10.85)
      -- (0.75,10.65) -- (0.75,10.35) -- cycle;
  \node[wf btn] (m1) at (2.65,10.55) {Számításelmélet I $\vee$};
  \node[wf btn] at (6.45,10.55) {Számításelmélet II $\vee$};
  \node[wf mark] at (9.5,10.55) {1};

  \node[wf box, anchor=west, font=\sffamily\scriptsize] at (0.6,9.3)
      {SZÁMÍTÁSELMÉLET ALAPJAI I--II};
  \node[anchor=west, font=\sffamily\large\bfseries] at (0.5,8.3)
      {Interaktív vizualizációs platform};
  \foreach \y/\l in {7.3/8.6, 6.85/8.2, 6.4/8.8, 5.95/5.0}
    \draw[wf line] (0.65,\y) -- (\l,\y);
  \foreach \y/\l in {5.1/8.4, 4.65/6.2}
    \draw[wf line] (0.65,\y) -- (\l,\y);
  \node[wf mark] at (9.5,8.3) {2};

  \draw[wf bar] (0,0) rectangle (10,1.3);
  \node[anchor=west, align=left] at (0.5,0.65)
      {Készítette: Szabó Anna\\ ELTE Informatikai Kar $\cdot$ 2026};
  \node[wf mark] at (9.5,0.65) {3};

  \begin{scope}[shift={(10.6,4.4)}]
    \node[anchor=south west, font=\sffamily\footnotesize\itshape] at (0,5.75)
        {Lenyíló témakörlista};
    \draw[wf frame, rounded corners=3pt] (0,0) rectangle (4.4,5.7);
    \fill[black!10, rounded corners=2pt] (0.1,2.1) rectangle (4.3,2.6);
    \foreach \t [count=\i] in {
        {Alapok 1: ábécé, szavak, \ldots},
        {Alapok 2: nyelvek és \ldots},
        Generatív grammatikák,
        Chomsky-hierarchia,
        Lineáris grammatikák \ldots,
        Determinisztikus véges a.,
        Nemdeterminisztikus v.\ a.,
        Reguláris kifejezések,
        \vdots}
      \node[anchor=west, font=\sffamily\scriptsize] at (0.2,5.95-0.6*\i) {\t};
    \node[wf mark] at (4.4,5.7) {4};
  \end{scope}
  \draw[densely dashed, black!50, ->] (m1.south east) to[out=-20, in=180] (10.6,9.6);
\end{tikzpicture}
}
    \caption{A főoldal vázlata}
    \label{fig:homepage}
\end{figure}

\clearpage

Egy témakör oldala (\ref{fig:module}.~ábra) felülről lefelé haladva
épül fel: (1) a helyet mutató morzsamenü és a cím, (2) az elméleti
összefoglaló tartalomblokkjai, (3) a \emph{Próbáld ki!} szakasz az
interaktív modullal, itt a bemeneti mezővel, a lejátszásvezérlővel és
az automata ábrájával, végül (4) a \emph{Teszteld magad!} szakasz a
kvízkérdésekkel. A (3) és (4) szakasz csak akkor jelenik meg, ha a
témakörhöz tartozik modul, illetve kvíz.

\begin{figure}[H]
    \centering
    \resizebox{!}{0.8\textheight}{\begin{tikzpicture}[wf]
  \draw[wf frame] (0,0) rectangle (13,20);

  \draw[wf bar] (0,19.1) rectangle (13,20);
  \draw[black!70, line width=0.7pt] (0.35,19.35) -- (0.35,19.65) -- (0.55,19.85)
      -- (0.75,19.65) -- (0.75,19.35) -- cycle;
  \node[wf btn] at (2.65,19.55) {Számításelmélet I $\vee$};
  \node[wf btn] at (6.45,19.55) {Számításelmélet II $\vee$};

  \node[anchor=west, font=\sffamily\scriptsize\itshape\bfseries, black!60] at (1,18.5)
      {Számításelmélet I > Reguláris nyelvek > Determinisztikus véges automata};
  \node[wf mark] at (12.5,18.5) {1};
  \node[font=\sffamily\Large\bfseries] at (6.5,17.6) {Determinisztikus véges automata};

  \draw[wf box] (1,14.6) rectangle (12,16.9);
  \node[anchor=north west, font=\sffamily\footnotesize\bfseries] at (1.15,16.85) {Definíció};
  \foreach \y/\l in {15.95/11.4, 15.5/10.8, 15.05/7.5}
    \draw[wf line] (1.4,\y) -- (\l,\y);
  \draw[wf box] (1,12.6) rectangle (12,14.3);
  \node[anchor=north west, font=\sffamily\footnotesize\bfseries] at (1.15,14.25) {Tétel};
  \foreach \y/\l in {13.4/11.2, 12.95/8.3}
    \draw[wf line] (1.4,\y) -- (\l,\y);
  \node at (6.5,12.2) {$\delta(q, a) = q'$};
  \node[anchor=west, font=\sffamily\footnotesize\itshape, black!60] at (1,11.75) {\ldots};
  \node[wf mark] at (12.5,16.9) {2};

  \node[anchor=west, font=\sffamily\large\bfseries] at (1,11.0) {Próbáld ki!};
  \node[wf mark] at (12.5,11.0) {3};
  \node[anchor=west] at (1,10.25) {Bemeneti szó:};
  \node[wf input, minimum width=3.6cm] at (3.4,10.25) {\textcolor{black!40}{$\varepsilon$ (üres szó)}};
  \foreach \s [count=\i] in {{$|\!\leftarrow$}, {$\leftarrow$}, {$\triangleright$}, {$\rightarrow$}, {$\rightarrow\!|$}}
    \node[wf btn, minimum width=0.7cm] at (7.4+0.8*\i,10.25) {\s};
  \draw[black!50, line width=1.5pt] (8.0,9.85) -- (12,9.85);
  \fill[black!70] (9.2,9.85) circle (0.09);

  \draw[wf box] (1,6.6) rectangle (12,9.6);
  \node[wf state, fill=black!20, line width=1.2pt] (q0) at (3,8.1) {$q_0$};
  \node[wf state] (q1) at (6.5,8.1) {$q_1$};
  \node[wf state, double, double distance=1.2pt] (q2) at (10,8.1) {$q_2$};
  \draw[->, >=stealth] (1.8,8.1) -- (q0);
  \draw[->, >=stealth] (q0) -- node[above] {$a$} (q1);
  \draw[->, >=stealth] (q1) -- node[above] {$b$} (q2);
  \draw[->, >=stealth] (q0) to[loop below] node[right] {$b$} (q0);
  \draw[->, >=stealth] (q2) to[out=-140, in=-40] node[below] {$a, b$} (q1);
  \node[anchor=west] at (1,6.1) {Konfigurációk: $(q_0, ab) \vdash (q_1, b) \vdash \ldots$};

  \node[anchor=west, font=\sffamily\large\bfseries] at (1,5.0) {Teszteld magad!};
  \node[wf mark] at (12.5,5.0) {4};
  \node[anchor=west, font=\sffamily\footnotesize\bfseries] at (1,4.45) {Könnyű};
  \draw[wf box] (1,1.55) rectangle (12,4.1);
  \node[anchor=west] at (1.3,3.7) {\textbf{1.} Elfogadja-e az automata az \emph{abba} szót?};
  \draw[draw=black!50, fill=black!15, rounded corners=2pt] (1.3,2.75) rectangle (11.7,3.25);
  \node[anchor=west] at (1.45,3.0) {\textbf{A)} Igen};
  \draw[draw=black!50, rounded corners=2pt] (1.3,2.1) rectangle (11.7,2.6);
  \node[anchor=west] at (1.45,2.35) {\textbf{B)} Nem};
  \node[anchor=west, font=\sffamily\scriptsize] at (1.3,1.8) {\textbf{Helyes!} Magyarázat \ldots};

  \draw[wf bar] (0,0) rectangle (13,1.2);
  \node[anchor=west, align=left] at (0.5,0.6)
      {Készítette: Szabó Anna\\ ELTE Informatikai Kar $\cdot$ 2026};
\end{tikzpicture}
}
    \caption{Egy témakör oldalának vázlata}
    \label{fig:module}
\end{figure}

\subsubsection{Adatformátum: a tartalommodell}
\label{sec:data}

A tervezés fő szempontjai a következők voltak:

\begin{itemize}
    \item a tartalom legyen független a megjelenítéstől a UI-tól;
    \item a tananyag legyen könnyen bővíthető;
    \item a teljes tananyag egységes formátumban tárolódjon;
    \item az interaktív modulok (DVA, NDVA, veremautomata stb.)
          ráépíthetők legyenek.
\end{itemize}

A tananyag egymásba ágyazott egységekből álljon: a témakörök
csoportokba tartoznak, minden témakör elméleti anyaga tartalomblokkok
sorozatából épül fel, tartozhatnak hozzá kvízkérdések, illetve interaktív
modul. A felhasználó haladása ettől függetlenül, a böngészőben
tárolódik. Az egyes egységeket a~\ref{tab:data-logical}.~táblázat
mutatja be.

\begin{table}[H]
    \centering
    \begin{tabularx}{\textwidth}{|l|>{\raggedright\arraybackslash}X|>{\raggedright\arraybackslash}X|}
        \hline
        \textbf{Egység} & \textbf{Mit tartalmaz} & \textbf{Példa} \\
        \hline
        Témakörcsoport & a tárgy egy nagyobb része
                       & Reguláris nyelvek \\ \hline
        Témakör        & cím, rövid összefoglaló, kapcsolódó előadások,
                         kapcsolódó interaktív modul
                       & Determinisztikus véges automaták, a hozzá
                         tartozó DVA-szimulátorral \\ \hline
        Tartalomblokk  & magának az elméleti anyag egy eleme (definíció, tétel,
                         bizonyítás, példa, képlet, táblázat \ldots)
                       & a DVA definíciója \\ \hline
        Kvízkérdés     & kérdés, válaszlehetőségek, helyes válasz,
                         magyarázat
                       & „Elfogadja-e az automata az \emph{abba} szót?” \\ \hline
        Haladás        & a megoldott feladatok és a beállítások, megtekintett témakörök
                       & melyik kvíz lett már kitöltve \\
        \hline
    \end{tabularx}
    \caption{A tananyag felépítése}
    \label{tab:data-logical}
\end{table}


%----------------------------------------------------------------------
\subsection{A követelmények validációja}
%----------------------------------------------------------------------
% ea02, 40. dia: prototípus + tesztesetek
% - a fejlesztés iteratív prototípusokban haladt
% - a felhasználói történetekből lesznek a funkcionális tesztesetek -> Tesztelés fejezet
A fejezetben kifejtett követelmények validálásának célja annak ellenőrzése,
 hogy a megfogalmazott követelmények megvalósíthatók, egymással összhangban
  állnak-e, valamint lefedik a rendszer célkitűzéseit.
   A követelmények összevetésre kerültek a projekt céljaival és a 
   meghatározott hatókörökkel. 
   Ennek eredményeként minden célkitűzéshez kapcsolódik legalább 
   egy funkcionális követelmény, továbbá nem szerepel olyan követelmény, 
   amely ellentmondana a hatókörön kívül eső funkcióknak. 
   A követelmények teljesülésének ellenőrzése a fejlesztés során 
   egységtesztek, manuális tesztesetek és a rendszer gyakorlati 
   kipróbálása segítségével történik.
    A részletes tesztelési eredményeket a Tesztelés fejezet ismerteti később.
    %======================================================================
\section{Tervezés}
\iffalse
%======================================================================
% TODO: ide megy át a thesis.tex "Áttekintés" (3 csomag + modulregiszter)
% és a "Rendszerarchitektúra" fejezet

\subsection{A tartalommodell}
\label{sec:content-design}

A blokkalapú modell előnye, hogy minden témakör ugyanazzal a
megjelenítő komponenssel rajzolható ki, új tartalom felvétele egyszerű,
és egy későbbi export (például PDF vagy nyomtatható nézet) is könnyen
megoldható.

A tananyagot TypeScript-modulokban tárolom. A szerkezet JSON-kompatibilis
(csak objektumokat, tömböket, szövegeket és számokat használ), de a
TypeScript-típusok miatt a fordító már szerkesztés közben jelzi, ha
egy blokk hiányos vagy hibás.

\subsubsection{Többnyelvű szöveg}

Minden megjelenő szöveg \texttt{LocalizedText} típusú. A magyar
változat kötelező, az angol opcionális, így a fordítás a meglévő
tartalom módosítása nélkül, fokozatosan pótolható.

\begin{lstlisting}
type LocalizedText = {
  hu: string;
  en?: string;
};
\end{lstlisting}

\subsubsection{Tartalomblokkok}

Egy témakör elméleti anyaga \texttt{ContentBlock} elemek sorozata.
A \texttt{ContentBlock} ún.\ diszkriminált unió (\emph{discriminated
union}): minden blokknak van egy \texttt{type} mezője, amely alapján a
fordító és a megjelenítő is el tudja dönteni, milyen további mezői
vannak. A jelenleg definiált blokktípusokat
a~\ref{tab:blocks}.~táblázat mutatja.

\begin{table}[H]
    \centering
    \begin{tabularx}{\textwidth}{|l|p{4.6cm}|X|}
        \hline
        \textbf{\texttt{type}} & \textbf{Mezők} & \textbf{Szerep} \\
        \hline
        \texttt{heading}    & \texttt{level} (1--3), \texttt{text}
                            & Fejezet- vagy szakaszcím \\ \hline
        \texttt{paragraph}  & \texttt{text}
                            & Folyó szöveg \\ \hline
        \texttt{definition} & \texttt{term}, \texttt{content}
                            & Fogalom definíciója \\ \hline
        \texttt{theorem}    & \texttt{kind}, \texttt{name?},
                              \texttt{statement}, \texttt{proof?}
                            & Tétel, lemma, állítás vagy következmény \\ \hline
        \texttt{proof}      & \texttt{title?}, \texttt{content}
                            & Bizonyítás, maga is blokkok sorozata \\ \hline
        \texttt{example}    & \texttt{title?}, \texttt{content}
                            & Kidolgozott példa \\ \hline
        \texttt{formula}    & \texttt{latex}
                            & Kiemelt képlet \\ \hline
        \texttt{list}       & \texttt{ordered}, \texttt{items}
                            & Számozott vagy felsorolásos lista \\ \hline
        \texttt{table}      & \texttt{headers}, \texttt{rows}
                            & Táblázat (pl.\ átmenettáblázat) \\ \hline
        \texttt{warning}    & \texttt{text}
                            & Figyelmeztetés tipikus hibára \\ \hline
        \texttt{funfact}    & \texttt{text}
                            & Érdekesség, kiegészítő megjegyzés \\
        \hline
    \end{tabularx}
    \caption{A tartalomblokkok típusai (a \texttt{?} opcionális mezőt jelöl)}
    \label{tab:blocks}
\end{table}

A \texttt{theorem} blokk \texttt{kind} mezője a
\texttt{tetel}, \texttt{lemma}, \texttt{allitas} és
\texttt{kovetkezmeny} értékek egyikét veheti fel. Egy definíció például
így néz ki:

\begin{lstlisting}
{
  type: "definition",
  term: { hu: "Determinisztikus véges automata" },
  content: { hu: "Az $A = \\langle Q, \\Sigma, \\delta, q_0, F \\rangle$ rendezett ötös, ahol ..." },
}
\end{lstlisting}

\subsubsection{Matematikai jelölések}

A tárgy anyaga sok matematikai jelölést tartalmaz (automaták
definíciói, grammatikák, pumpálási lemmák, a CYK algoritmus stb.),
ezért a képleteket egységesen \LaTeX{} szintaxissal tárolom:

\begin{itemize}
    \item szövegen belüli képlet: \verb|$...$| a szövegmezőben;
    \item kiemelt képlet: külön \texttt{formula} blokk;
    \item a sztringekben a \verb|\| karaktert duplán kell írni
          (\verb|\\|).
\end{itemize}

\subsubsection{Témakörök metaadatai}

A tartalom és a navigáció szétválasztása érdekében minden témakörhöz
külön metaadat (\texttt{TopicMeta}) tartozik. Ez szolgál a
navigációhoz, a sorrendezéshez, a csoportosításhoz, az előadásokhoz
kötéshez, és ezen keresztül kapcsolódik a témakörhöz az
interaktív modul. Egy témakör (\texttt{Topic}) a metaadatokból, a
tartalomblokkok listájából és opcionálisan kvízkérdésekből áll.

\begin{table}[H]
    \centering
    \begin{tabularx}{\textwidth}{|l|l|X|}
        \hline
        \textbf{Mező} & \textbf{Típus} & \textbf{Jelentés} \\
        \hline
        \texttt{id}       & \texttt{string}        & Egyedi azonosító, az URL-ben is ez szerepel \\ \hline
        \texttt{order}    & \texttt{number}        & Megjelenési sorrend \\ \hline
        \texttt{group}    & \texttt{TopicGroup}    & Témakörcsoport \\ \hline
        \texttt{title}    & \texttt{LocalizedText} & Cím \\ \hline
        \texttt{summary}  & \texttt{LocalizedText} & Rövid összefoglaló \\ \hline
        \texttt{lectures} & \texttt{number[]}      & Kapcsolódó előadások sorszáma \\ \hline
        \texttt{module}   & \texttt{string | null} & A kapcsolódó interaktív modul azonosítója \\
        \hline
    \end{tabularx}
    \caption{A \texttt{TopicMeta} mezői}
    \label{tab:topicmeta}
\end{table}

A témakörök a tárgy logikai felosztását követve öt csoportba
(\texttt{TopicGroup}) sorolódnak:

\begin{table}[H]
    \centering
    \begin{tabular}{|l|l|}
        \hline
        \textbf{Azonosító} & \textbf{Jelentés} \\
        \hline
        \texttt{foundations}       & Alapfogalmak \\ \hline
        \texttt{regular}           & Reguláris nyelvek \\ \hline
        \texttt{context-free}      & Környezetfüggetlen nyelvek \\ \hline
        \texttt{context-sensitive} & Környezetfüggő és 0-típusú nyelvek \\ \hline
        \texttt{computability}     & Kiszámíthatóság (II. rész) \\
        \hline
    \end{tabular}
    \caption{Témakörcsoportok}
    \label{tab:groups}
\end{table}

% TODO: témakörök száma és listája a végén (most 16 mappa a topics/ alatt)  \fi